\subsection{Scheduling Policy and Goals}
Both FreeBSD and NetBSD are direct descendants of the Berkeley Software Distribution (BSD) Unix lineage, and both handle processor management through kernel-level schedulers. But while they share common roots, they have each developed quite different approaches to how they prioritize tasks, manage threads, and decide which process gets CPU time. Understanding these differences gives us a clearer picture of what each OS was designed to do well.

FreeBSD has two schedulers available in its kernel. The default one is called the ULE scheduler, found in /sys/kern/sched\_ule.c. There is also the older traditional 4.4BSD scheduler, which is still maintained but no longer used by default \parencite{watson2015}. The key thing about FreeBSD's scheduling goals is that it is trying to do several things at once: it wants to be fair to all threads, responsive to interactive tasks like text editors and web browsers, and efficient for batch or background work.
The ULE scheduler organizes threads into three broad categories: interactive threads, non-interactive (batch) threads, and real-time threads. It assigns priorities based on behavior, specifically by looking at how much time a thread voluntarily sleeps versus how much time it actually runs. A thread that sleeps a lot and wakes up in bursts, like one waiting for keyboard input, is tagged as interactive and given a higher scheduling priority. A thread that just constantly consumes CPU without releasing it gets treated as a batch job and is pushed further down the priority order \parencite{watson2015}.
FreeBSD defines priority numerically from 0 to 255, where a lower number means higher priority. The scheduling classes are organized by range: interrupt threads get priority 0 to 47, real-time threads get 48 to 79, top-half kernel threads get 80 to 119, timeshare threads get 120 to 223, and idle threads get 224 to 255 \parencite{watson2015}. The main goal is to "balance resource utilization against the time that it takes for a program to complete," as described in the FreeBSD Internals book \parencite{watson2015}

NetBSD's scheduling philosophy comes from a different angle. NetBSD has historically prioritized portability and correctness over any one platform's performance characteristics. The scheduler in NetBSD is also derived from the traditional BSD scheduler, but with an emphasis on being usable across a very wide range of hardware platforms, from embedded devices to servers. NetBSD's scheduling priorities are set and controlled through splbio(), spltty(), and similar interrupt priority level (IPL) mechanisms that manage processor priority at a hardware level (Silvers, 2002).
NetBSD also relies on a priority-based preemptive scheduling model where the kernel decides which thread to run based on its current priority. The scheduler aims to give fair CPU time to all running processes, with mechanisms to ensure that I/O-bound processes are favored over CPU-bound ones since I/O operations tend to complete quickly and the process quickly goes back to waiting. In the NetBSD source tree, the scheduler logic is located in the kern/ directory alongside other core kernel management code, covering process and kernel thread management as a unified subsystem \parencite{feyrer2002}.
In terms of comparing goals: FreeBSD places a stronger emphasis on interactive responsiveness and has engineered a more sophisticated interactivity detection system to support that. NetBSD leans toward broad correctness and portability. Both systems use dynamic priority adjustment, but FreeBSD's mechanisms are considerably more elaborate, especially with the ULE scheduler's interactivity scoring algorithm.

\subsection{Scheduling Schemes}
FreeBSD uses a multilevel feedback queue system that has been refined significantly over the years, particularly through the development of ULE. The scheduler is split into two parts: a low-level scheduler and a high-level scheduler. The low-level scheduler is responsible for the actual decision of which thread to run next and is called every time a thread blocks. For efficiency, it just picks the first thread from the highest-priority non-empty run queue for the current CPU. The high-level scheduler handles the harder work of assigning priorities and deciding which CPU's run queue a thread should be placed in \parencite{watson2015}.
The data structure that holds runnable threads is an array of 64 run queues per CPU, where a thread's priority divided by four determines which queue it lands in. There is also a bit vector called rq\_status associated with this array that tracks which queues are non-empty, allowing the scheduler to quickly find the highest priority runnable thread without scanning all 64 queues one by one \parencite{watson2015}.
CPU time is allocated using time slices. For interactive threads, the ULE scheduler places them in the realtime queue and schedules them like real-time round-robin threads. Batch threads go into a timeshare calendar queue, which is managed differently. Rather than ordering threads strictly by priority, the calendar queue rotates through positions over time, so even the lowest-priority batch thread will eventually get CPU time. This ensures that no timeshare thread starves completely \parencite{watson2015}. The time slice length itself is adjusted dynamically based on system load, with a minimum size to avoid excessive context-switching overhead.
All of this is fully preemptive. An interrupt or higher-priority thread can preempt anything that is currently running if it has a higher priority. The kernel uses an Asynchronous System Trap (AST) mechanism to force context switches when needed, even mid-user-mode execution \parencite{watson2015}.

Though NetBSD also uses a priority-based preemptive scheduling model that traces back to the traditional BSD scheduler, the kernel keeps a set of run queues where runnable threads are placed according to their priority. When the scheduler needs to pick the next thread to run, it selects the thread from the highest non-empty priority queue. Within a queue, threads are run using round-robin scheduling so that threads at the same priority level each get a fair turn (Silvers, 2002).
The scheduling decisions in NetBSD are also preemptive, meaning a higher-priority thread can interrupt a lower-priority one. The interrupt priority level system, managed through functions like splbio() and spltty(), controls when the processor can be interrupted, and these levels are carefully managed to avoid deadlocks between interrupt handlers and regular thread execution \parencite{netBSDnd}.
One key difference from FreeBSD is that NetBSD's scheduling structure is more straightforward. FreeBSD's ULE scheduler has significantly more moving parts, including the interactivity scoring, the CPU topology tree for multiprocessor decisions, and the adaptive idle feature. NetBSD's approach is simpler and arguably more predictable, which fits well with its goal of running reliably on a wide variety of hardware.

\subsection{Context Switching}
Context switching in FreeBSD is divided into two types: voluntary and involuntary. A voluntary context switch happens when a thread calls the sleep() routine and gives up the CPU because it needs to wait for something, like I/O completion or a lock to become available. An involuntary context switch happens when a thread uses up its time quantum or when the system detects a higher-priority thread that should run instead (McKusick et al., 2015, p. 99).
The actual mechanism is split between machine-independent code in mi\_switch() and machine-dependent code in cpu\_switch(), which is typically written in assembly language for efficiency. The context saved and restored includes the Thread State Block (TSB), which holds the general-purpose registers, stack pointers, program counter, processor-status word, and memory-management registers. The user-mode execution state is stored on the kernel execution stack when a thread is in kernel mode, and the kernel makes sure this state is always properly saved whenever the kernel is entered (McKusick et al., 2015, p. 100-101).
The cost of context switching in FreeBSD is kept reasonably low through several design decisions. Each thread has its own pre-allocated kernel stack, so there is no memory allocation cost during a switch. The turnstile and sleepqueue data structures are also pre-allocated per thread so that blocking on a lock or sleeping for an event does not require dynamic memory allocation either \parencite{watson2015}. The tradeoff is that frequent context switches mean more time spent in the scheduler and more cache pollution, which is why the ULE scheduler dynamically adjusts time slice lengths to keep the balance between responsiveness and throughput.
FreeBSD supports both SMP (Symmetric Multi-Processing) and NUMA (Non-Uniform Memory Access) architectures, which adds complexity to context switching because the kernel needs to manage per-CPU run queues, locking, and cache considerations during thread migration between CPUs \parencite{watson2015}.

In NetBSD, context switching follows the same fundamental BSD model. When the currently running thread needs to be suspended, the kernel saves its state and selects the next runnable thread. The preempt() function is called in critical sections of the NetBSD kernel when a higher-priority thread needs to run, as seen in the kernel source code \parencite{netBSDnd}.
The cost of context switching in NetBSD is comparable to FreeBSD at a basic level, since both share the same BSD roots. However, because NetBSD targets a broad range of hardware including embedded and resource-constrained systems, the implementation tends to be leaner in some aspects. NetBSD relies on the uvm\_scheduler() as an infinite loop that manages scheduling at the virtual memory level as well, tying together memory management and process scheduling \parencite{netBSDnd}.
The main tradeoff in NetBSD's context switching is similar to FreeBSD's: too-frequent switches hurt throughput because of cache thrashing and scheduler overhead, while infrequent switches hurt interactivity. NetBSD's simpler scheduler means it has fewer heuristics to manage this tradeoff, which can be a disadvantage on desktop or heavily interactive workloads but is less of a concern on server or embedded platforms where workloads tend to be more predictable.

\subsection{Process and Thread Handling}
FreeBSD distinguishes clearly between processes and threads. A process is basically a container that holds an address space and system resources. Within each process, there can be one or many threads, where each thread is the actual unit of execution. Threads within the same process share the address space and other resources but are scheduled independently of each other \parencite{watson2015}.
FreeBSD uses a 1:1 threading model, meaning every user-level thread is backed by a kernel-level thread. This is implemented through the POSIX threading API (pthreads), specifically following the Pthreads model. Historically, FreeBSD went through a period where it used an N:M model, where many user-level threads were supported by fewer kernel threads, but this was found to be problematic for modern applications that regularly enter the kernel. The 1:1 model offers better performance for these applications because kernel operations on one thread do not block other threads in the same process \parencite{watson2015}.
Each thread in FreeBSD has its own thread structure that holds scheduling information, a kernel stack, the TSB, and machine-dependent state. The process structure itself holds everything that is shared across the thread group, including the address space, file descriptors, signal actions, credentials, and resource limits \parencite{watson2015}.
The advantage of the 1:1 model is that each thread can independently make system calls, block on I/O, and be scheduled on any available CPU. The disadvantage is that creating and managing large numbers of threads has more kernel overhead compared to a user-level thread model. However, for most real-world applications, the 1:1 model's benefits clearly outweigh this cost.

NetBSD has a similar process-and-thread architecture, but with an interesting history. NetBSD historically used kernel threads for internal kernel operations, and the user-land thread model evolved differently from FreeBSD. Like FreeBSD, NetBSD now supports kernel-level threads through a 1:1 model for POSIX threads, where each user thread corresponds to a kernel thread. The kernel creates various kernel threads at startup for tasks like page daemon management, filesystem syncing, and the process reaper, all managed through the kthread\_create1() interface \parencite{netBSDnd}.
The main process structure in NetBSD is struct proc, which contains per-process information including emulation data, credentials, signal state, and timing. This maps closely to the FreeBSD model, both systems having inherited this structure from the same BSD lineage \parencite{netBSDnd}. Threads are scheduled independently just as in FreeBSD, and they can be preempted by higher-priority threads.
One area where NetBSD is somewhat different is in its broader concept of lightweight processes and its support for emulating other operating systems through compat layers in the kernel. The compat subsystem allows NetBSD to run binaries from Linux, FreeBSD, SunOS, and others, which adds complexity to how processes are created and managed. This is a unique feature of NetBSD's design that reflects its focus on portability \parencite{feyrer2002}.

\subsection{Behavior under different workloads}
For CPU-intensive work, like compiling a large program or running a mathematical computation, both FreeBSD and NetBSD handle it reasonably well. The key difference is in how gracefully they manage such workloads alongside other tasks. In FreeBSD, a CPU-heavy process will quickly accumulate run time and have its priority lowered, moving it into the batch category in ULE's scheduling model. This means that interactive tasks like a running terminal or a text editor will still get reasonable priority and will not feel sluggish \parencite{watson2015}.
There is even a specific mechanism in FreeBSD where the scheduling priority of a child process is propagated back to the parent. So if a build system like make starts many CPU-intensive compilation jobs, the make process itself has its priority gradually lowered as its children consume CPU. New compilation steps then start at this lower priority, which prevents the compilation from monopolizing the system at the expense of interactive applications \parencite{watson2015}.
NetBSD handles CPU-bound workloads similarly from a priority standpoint, but without the same level of interactivity scoring sophistication that FreeBSD has. On a single-user system running a CPU-intensive task, NetBSD should perform comparably, but FreeBSD will likely feel more responsive to the interactive user because of its more aggressive interactivity detection.

I/O-bound workloads are where both systems tend to perform well. A thread that voluntarily sleeps while waiting for disk or network I/O will have its priority raised when it wakes up, because both BSD-derived schedulers recognize that sleeping threads should get a quick turn when they are ready. In FreeBSD, threads that go to sleep in the kernel get a higher priority when they are awakened because they typically hold shared kernel resources, and the system wants them to use those resources and release them quickly before another thread gets blocked waiting \parencite{watson2015}.
For server workloads with lots of concurrent I/O operations, FreeBSD is particularly well known for its performance and stability. The infoworld.com assessment of FreeBSD noted its exceptional stability on servers handling large volumes of I/O, with one FreeBSD server running over 1,000 days uptime under heavy email and DNS load without issue (Venezia, 2011). This is partly a testament to both the scheduler's I/O-handling design and the overall system architecture.

This is where FreeBSD clearly has an edge. The ULE scheduler's interactivity score system is specifically designed to detect programs that behave like interactive applications, those that spend most of their time sleeping and wake up briefly to respond to user input. These programs are kept in a high-priority class that runs alongside real-time threads, so they consistently get quick CPU access \parencite{watson2015}.
FreeBSD also handles complex GUI applications like web browsers, which can have mixed behavior, with short bursts of heavy computation mixed with long idle periods. The ULE scheduler keeps a moving average of sleep/run history and decays it over time, which means a browser does not get permanently downgraded to batch status just because it had a brief period of CPU-intensive activity \parencite{watson2015}.
NetBSD, while capable, is less optimized for interactive desktop use. It is more commonly deployed in server or embedded contexts where interactive responsiveness is less of a concern. For user-facing applications, FreeBSD provides a noticeably better experience.

\subsection{Multicore and Parallel Processing}
Multiprocessor support was one of the primary motivations for developing the ULE scheduler in the first place. The traditional 4.4BSD scheduler had a single global run queue protected by a single global lock, which became a severe bottleneck on multiprocessor systems. ULE addresses this by maintaining per-CPU run queues, so each CPU has its own set of realtime, timeshare, and idle queues. This means multiple CPUs can each independently select their next thread to run without contending for a shared lock \parencite{watson2015}.
FreeBSD's multiprocessor scheduling uses a CPU topology tree to model the relationships between CPUs. The scheduler understands which CPUs share caches, which are on the same processor chip, and which are on different physical packages. When deciding where to run a thread, it prefers to run it on a CPU that shares a cache with the CPU where it last ran, since this reduces cache miss penalties from migrating the thread's working set to a new cache. The closer the CPUs are in the topology tree, the cheaper the migration \parencite{watson2015}.
Load balancing across cores happens through three mechanisms: when a thread becomes runnable (it can be pushed to an idle CPU immediately), when a CPU goes idle (it steals work from other CPUs), and through a periodic long-term load balancer that runs roughly once per second to fix longer-term imbalances. The idle CPU checks a shared bitmask of idle CPUs, which is much cheaper than scanning all run queues \parencite{watson2015}.
The tradeoff between cache affinity and load balancing is explicitly managed in ULE. When a thread wakes up, the scheduler first tries to place it on the same CPU where it previously ran, as long as the time it was asleep was short enough that its cache data is likely still warm. If the sleep was too long, or if affinity cannot be maintained without leaving CPUs idle, it will migrate the thread \parencite{watson2015}. FreeBSD also supports hard CPU affinity through the cpuset interface, letting applications or administrators explicitly pin processes to specific CPUs

NetBSD supports SMP as well, and has evolved over time to improve multiprocessor performance. The NetBSD kernel structure includes per-CPU scheduling state (stored in the ci\_schedstate field of the CPU info structure), which reflects its effort to manage per-CPU scheduling separately. Preemption is handled through calls like preempt() in the kernel when a higher-priority thread becomes runnable \parencite{netBSDnd}.
NetBSD's multiprocessor support is functional but generally considered less sophisticated than FreeBSD's in terms of the topology-aware load balancing and adaptive idle features that ULE provides. NetBSD's portability focus means it has to work correctly on a very wide variety of hardware, including older or exotic multiprocessor systems, which shapes its implementation choices. The performance on modern multi-core x86 servers is solid, but for workloads that heavily benefit from NUMA-aware scheduling or advanced cache affinity, FreeBSD's ULE has more mature tooling.
Both systems support CPU affinity for threads, but FreeBSD's implementation is more refined and better documented. FreeBSD's design scales well from single CPUs to 512-thread network processors according to its own developers \parencite{watson2015}.

\subsection{Starvation, Fairness and Priority Handling}
FreeBSD takes starvation prevention seriously, and the ULE scheduler has a specific mechanism to address it for batch threads. The timeshare queue in ULE is managed as a calendar queue. Rather than simple priority ordering, the calendar queue rotates over time, and a thread's insertion position in the queue is based on its relative priority. A lower-priority thread is placed farther from the current position, but the queue pointer advances over time, meaning even the lowest-priority batch thread will eventually reach the front and get to run. The scheduler is designed so that no timeshare thread can be completely starved out (McKusick et al., 2015, p. 121-122).
Priority adjustment over time happens through two mechanisms. For interactive threads, the interactivity score is continuously recalculated based on recent sleep and run history, decaying the history gradually so that older behavior has less influence. This allows a temporarily CPU-intensive thread to recover its interactive status once it goes back to normal behavior. For batch threads, priority is determined by CPU utilization estimates that also decay over time, preventing a thread that has used a lot of CPU in the past from being permanently stuck at low priority \parencite{watson2015}.
FreeBSD also implements priority inheritance through turnstile data structures. If a lower-priority thread is holding a lock that a higher-priority thread needs, the lower-priority thread temporarily gets its priority boosted to match the waiting thread. This prevents priority inversion, where a high-priority thread is effectively blocked by a low-priority one (McKusick et al., 2015, p. 103). This is a form of priority propagation rather than priority aging, but it serves a similar function of preventing unfair blocking.

NetBSD uses a priority-based scheduler where priorities are adjusted dynamically, similar to the traditional BSD approach. Threads that do more I/O will tend to accumulate less CPU time and therefore maintain a higher priority, while CPU-heavy threads see their priority gradually lowered. This is the same feedback mechanism found in the original BSD scheduler and forms the basis of starvation prevention: CPU-hungry processes get lower priority over time, giving other processes a chance to run.
However, NetBSD's starvation prevention mechanisms are somewhat less elaborate than FreeBSD's. The traditional BSD approach of periodically scanning all runnable threads to adjust their priorities (decay-based priority adjustment) is simpler but works reasonably well for most workloads. The main risk in any priority-based system is that if the priority adjustment does not happen frequently enough, some threads could experience significant delays. NetBSD's approach is conservative and well-tested, but it lacks the more sophisticated calendar queue mechanism that FreeBSD uses to provide stronger starvation guarantees for batch threads.
NetBSD uses the lock manager's priority field to control the priority at which a thread resumes after sleeping, which provides a basic form of priority control when threads are woken up from blocking operations \parencite{netBSDnd}. This is straightforward and works correctly but does not include the turnstile-based priority propagation that FreeBSD uses

\subsection{Summary}
FreeBSD and NetBSD share the same BSD scheduling roots but have evolved in noticeably different directions. FreeBSD's ULE scheduler is the more sophisticated of the two, featuring per-CPU run queues, a topology-aware load balancer, turnstile-based priority inheritance, and an interactivity scoring algorithm that dynamically keeps interactive threads prioritized over batch jobs. These features make FreeBSD the stronger choice for responsiveness, multicore scalability, and mixed workloads where interactive applications and background tasks compete for CPU time. Its calendar queue design also provides stronger starvation guarantees for low-priority threads, and the entire scheduler is built to run in constant time regardless of how many threads are in the system (McKusick et al., 2015, pp. 114-125).
NetBSD's scheduler is simpler and more conservative, prioritizing correctness and portability over advanced heuristics. It uses a straightforward priority-based preemptive model that works reliably across a very wide range of hardware, from embedded systems to servers. This makes it well-suited for environments where predictable behavior and broad platform support matter more than peak interactive performance. The tradeoff is that NetBSD offers less sophisticated interactivity detection, less mature multiprocessor load balancing, and lacks priority inheritance mechanisms like turnstiles. For server workloads with more uniform behavior, or for specialized deployments like network appliances and SAN arrays where NetBSD is commonly found (Venezia, 2011), these limitations matter less. But in a mixed-workload scenario with both interactive applications and background tasks, FreeBSD's scheduler gives users a noticeably smoother experience.